%% kapa: exam=1 round=33 year=2022 period=1 form=A subject=economics subject_ko=경제학원론 seq=15 number=55
%% source: Q-Net 감정평가사 기출문제 (cst00309) artlSeq=5209997 file=raw/1008/5209997/2022년도 제33회 감정평가사 1차시험 1교시(민법,경제학,부동산학).pdf
%% conversion: kordoc 4.24.51 -> tools/gen1.py (auto); review=proofread-sonnet; flags=none
\begin{question}{55}
두 재화 X재와 Y재를 소비하는 갑은 가격이 (P$_{\mathrm{X}}$, P$_{\mathrm{Y}}$) = (1, 4)일 때 소비조합 (X, Y) = (6, 3), 가격이 (P$_{\mathrm{X}}$, P$_{\mathrm{Y}}$) = (2, 3)으로 변화했을 때 소비조합 (X, Y) = (7, 2), 그리고 가격이 (P$_{\mathrm{X}}$, P$_{\mathrm{Y}}$) = (4, 2)으로 변화했을 때 소비조합 (X, Y) = (6, 4)을 선택하였다. 이에 관한 설명으로 옳은 것을 모두 고른 것은?
\begin{framedtext}
ㄱ. 소비조합 (X, Y) = (6, 3)이 소비조합 (X, Y) = (7, 2)보다 직접 현시선호되었다. \par
ㄴ. 소비조합 (X, Y) = (6, 4)이 소비조합 (X, Y) = (7, 2)보다 직접 현시선호되었다. \par
ㄷ. 소비조합 (X, Y) = (6, 3)이 소비조합 (X, Y) = (6, 4)보다 직접 현시선호되었다. \par
ㄹ. 선호체계는 현시선호이론의 약공리를 위배한다.
\end{framedtext}
\begin{choices}
  \item ㄱ, ㄴ
  \item ㄱ, ㄷ
  \item ㄱ, ㄹ
  \item ㄴ, ㄷ
  \item ㄷ, ㄹ
\end{choices}
\end{question}
\answer{1}
